% !TEX program = xelatex
\documentclass[UTF8,12pt,a4paper]{ctexart}

\usepackage[a4paper,margin=2.5cm]{geometry}
\usepackage{amsmath,amssymb,bm}
\usepackage{booktabs}
\usepackage{graphicx}
\usepackage{float}
\usepackage{xcolor}
\usepackage{listings}
\usepackage{hyperref}
\usepackage{enumitem}

\hypersetup{
    colorlinks=true,
    linkcolor=blue!55!black,
    citecolor=blue!55!black,
    urlcolor=blue!55!black
}

\lstset{
    language=C,
    basicstyle=\ttfamily\small,
    keywordstyle=\color{blue!70!black},
    commentstyle=\color{green!45!black},
    stringstyle=\color{red!60!black},
    numbers=left,
    numberstyle=\tiny\color{gray},
    frame=single,
    breaklines=true,
    columns=fullflexible,
    showstringspaces=false
}

\newcommand{\ReportTitle}{1D-Riemann 问题 MacCormack 格式： 方法与 Code }
\newcommand{\ReportAuthor}{邵桂博}
\newcommand{\ReportCourse}{计算流体力学(Project 2)}
\newcommand{\ReportDate}{\today}

\title{\ReportTitle}
\author{\ReportAuthor\\\ReportCourse}
\date{\ReportDate}

\begin{document}

\maketitle
\begin{abstract}
本文介绍我的1D Riemann问题-MacCormack离散格式的求解思路与代码实现思路；目录中1 问题描述与控制方程，内容是①概述待解的问题②明确求解的任务；具体的求解思路为 2 数值求解方法； 3 程序结构与代码逻辑（伪代码形式呈现）； 4 关键实现细节； 5 本部分小结
预测--校正格式、人工粘性处理以及 C 语言程序实现。

\textbf{Keywords：} Riemann 问题；MacCormack 格式；人工粘性
\end{abstract}

\tableofcontents
\newpage

\section{问题描述与控制方程}

\subsection{一维 Riemann 问题}

\subsubsection{此问题的坐标系与控制方程}
    对于一维问题，坐标系自然地取为以右为正方向的一维直线坐标系，
    在此坐标系下控制方程为
\begin{equation}
\left\{
\begin{aligned}
\frac{\partial \rho}{\partial t}
&+
\frac{\partial (\rho u)}{\partial x}
=0,\\[6pt]
\frac{\partial (\rho u)}{\partial t}
&+
\frac{\partial (\rho u^2+p)}{\partial x}
=0,\\[6pt]
\frac{\partial (\rho E)}{\partial t}
&+
\frac{\partial (\rho E u+p u)}{\partial x}
=0.
\end{aligned}
\right.
\end{equation}

其中，单位质量流体的总能量为
\begin{equation}
E=e+\frac{u^2}{2}
=
\frac{p}{\rho(\gamma-1)}
+
\frac{u^2}{2}.
\end{equation}

\subsubsection{左右初始状态}

\noindent 1. 激波间断--接触间断--膨胀波问题之 Sod 问题
\begin{equation}
\begin{cases}
x<0: & (\rho,u,p)_L=(1,\ 0,\ 1),\\
x\geq 0: & (\rho,u,p)_R=(0.125,\ 0,\ 0.1).
\end{cases}
\end{equation}

\noindent 2. 激波间断--接触间断--膨胀波问题之 Lax 问题
\begin{equation}
\begin{cases}
x<0: & (\rho,u,p)_L=(0.445,\ 0.698,\ 3.528),\\
x\geq 0: & (\rho,u,p)_R=(0.5,\ 0,\ 0.571).
\end{cases}
\end{equation}

\noindent 3. 双膨胀波问题之亚声速后退问题
\begin{equation}
\begin{cases}
x<0: & (\rho,u,p)_L=(1,\ -2,\ 4),\\
x\geq 0: & (\rho,u,p)_R=(1,\ 2,\ 4).
\end{cases}
\end{equation}

\noindent 4. 双膨胀波问题之 Sjogreen supersonic Expansion 问题
\begin{equation}
\begin{cases}
x<0: & (\rho,u,p)_L=(1,\ -2,\ 0.4),\\
x\geq 0: & (\rho,u,p)_R=(1,\ 2,\ 0.4).
\end{cases}
\end{equation}

\noindent 5. 接触间断--双膨胀波问题
\begin{equation}
\begin{cases}
x<0: & (\rho,u,p)_L=(1,\ -0.2,\ 0.5),\\
x\geq 0: & (\rho,u,p)_R=(0.5,\ 0.5,\ 0.5).
\end{cases}
\end{equation}

\noindent 6. 接触间断--双激波间断问题
\begin{equation}
\begin{cases}
x<0: & (\rho,u,p)_L=(0.4,\ 0.5,\ 1.0),\\
x\geq 0: & (\rho,u,p)_R=(1,\ -0.5,\ 0.9).
\end{cases}
\end{equation}

\noindent 7. 接触间断问题
\begin{equation}
\begin{cases}
x<0: & (\rho,u,p)_L=(10,\ 1.0,\ 2.0),\\
x\geq 0: & (\rho,u,p)_R=(1.0,\ 1.0,\ 2.0).
\end{cases}
\end{equation}

\subsubsection{计算目标}

本项目采用 MacCormack 格式求解一维 Riemann 问题的数值解，并将数值解与
半解析解（精确解）进行对比。具体研究目标如下。

\begin{enumerate}[leftmargin=2em]
    \item 求解 1.1.2 中的几个场初值问题，具体如下：
    \begin{itemize}
        \item 将 MacCormack 数值解与精确解进行对比；
        \item 比较加入和不加入人工粘性时的计算结果；
        \item 研究人工粘性经验参数 $\beta$ 对数值解的影响；
        \item 研究 CFL 数对数值解的影响，并讨论人工粘性系数与 CFL 数之间
        并非简单单调关系的现象；
        \item 以密度作为人工粘性开关函数，并进一步比较采用压力或速度作为
        开关函数时的差异；
        \item 在不加入人工粘性的情况下，通过加密网格考察色散效应是否得到改善，
        并分析 MacCormack 格式固有色散误差与网格分辨率之间的关系。
    \end{itemize}

    \item 对于 Sjogreen supersonic Expansion 超声速后退问题，重点考察真空区：
    \begin{itemize}
        \item MacCormack 格式能否稳定计算真空或近真空状态；
        \item 真空区内密度、压力和速度应如何处理；
        \item 数值保护、人工粘性和网格分辨率对结果的影响。
    \end{itemize}

    \item 对于接触间断--双膨胀波问题，分析接触间断和两侧膨胀波的捕捉效果。

    \item 对于接触间断--双激波问题，分析激波位置、间断厚度及数值振荡。

    \item 对于纯接触间断问题，考察 MacCormack 格式对接触间断的分辨能力，

    并在条件允许时与其他高阶或高分辨率格式进行比较。通过该算例说明不同数值格式
    各自擅长捕捉的波结构存在差异，不能仅依据格式阶数判断其在所有问题上的表现。
\end{enumerate}

\subsubsection{Project 2与Project 0 半解析解之间的关系}
Project 0 作为半解析解（尽管在求解超越方程时使用了数值计算方法），可以理解为一种Benchmark，检验同问题但利用 DNS 方法、采用 MacCormack 格式得到的解的精度、稳定性等性质。

我们的原始变量都为W（原始量），作为最终需要输出的解。
\begin{equation}
\bm{W}(x,0)=
\begin{cases}
(\rho_L,u_L,p_L)^{\mathrm T}, & x<x_0,\\
(\rho_R,u_R,p_R)^{\mathrm T}, & x>x_0.
\end{cases}
\end{equation}

\subsection{一维 Euler 方程}

\begin{equation}
\frac{\partial \bm{Q}}{\partial t}
+
\frac{\partial \bm{F}(\bm{Q})}{\partial x}
=0,
\end{equation}

\begin{equation}
\bm{Q}=
\begin{bmatrix}
\rho\\
\rho u\\
\rho E
\end{bmatrix},
\qquad
\bm{F}=
\begin{bmatrix}
\rho u\\
\rho u^2+p\\
u(\rho E+p)
\end{bmatrix}.
\end{equation}

总能量满足

\begin{equation}
E=\frac{p}{\rho(\gamma-1)}+\frac{u^2}{2}.
\end{equation}


\section{数值求解方法}

\subsection{守恒量与原始量 - 未知变量储存的策略}

我的程序同时保存守恒量
$\bm{Q}=[\rho,\rho u,\rho E]^{\mathrm T}$ 和原始量
$\bm{W}=[\rho,u,p]^{\mathrm T}$ 。其中守恒量是时间推进主状态，
原始量由守恒量反解，用于声速、CFL、人工粘性、输出和诊断。

原因讨论如下。

虽然本问题求解的是同一组守恒型方程，但根据推进过程中选取的核心变量不同，
可以采用两种内存储存方案。

\subsubsection{以守恒量 \texorpdfstring{$\bm{Q}$}{Q} 为核心变量}

第一种方案在每个网格点储存
\begin{equation}
\bm{Q}=
\begin{bmatrix}
\rho & \rho u & \rho E
\end{bmatrix}^{\mathrm T},
\end{equation}
并由数值格式直接更新守恒量。MacCormack、有限体积、Roe 和 HLLC
等守恒型数值方法天然以 $\bm{Q}$ 作为推进变量。

该方案的基本流程为
\begin{equation}
\text{内存分配}
\longrightarrow
\bm{Q}
\longrightarrow
\text{时间循环}
\longrightarrow
\text{更新 }\bm{Q}
\longrightarrow
\text{下一时间步}.
\end{equation}

通量、声速、CFL 条件和人工粘性系数通常需要使用原始量，因此每个时间步还需进行
\begin{equation}
\bm{Q}\longrightarrow\bm{W},
\qquad
\bm{W}\longrightarrow\bm{F},
\qquad
\bm{W}\longrightarrow\mu,
\end{equation}
其中 $\mu$ 表示人工粘性系数。控制方程的非线性主要集中在守恒量到原始量的转换、
矢通量构造和人工粘性系数计算中。

\subsubsection{以原始量 \texorpdfstring{$\bm{W}$}{W} 为核心变量}

第二种方案在每个网格点储存
\begin{equation}
\bm{W}=
\begin{bmatrix}
\rho & u & p
\end{bmatrix}^{\mathrm T},
\end{equation}
再由 $\bm{W}$ 组合出守恒量 $\bm{Q}(\bm{W})$、通量
$\bm{F}(\bm{W})$ 和人工粘性系数。其基本流程为
\begin{equation}
\text{内存分配}
\longrightarrow
\bm{W}
\longrightarrow
\bm{Q}(\bm{W}),\ \bm{F}(\bm{W}),\ \mu(\bm{W})
\longrightarrow
\text{更新}
\longrightarrow
\text{下一时间步}.
\end{equation}

Project 0 利用特征关系求解 Riemann 问题的半解析解时，主要采用原始量描述左右状态、
激波、膨胀波和接触间断，因此这种组织方式更符合其理论推导逻辑。

\subsubsection{本程序采用的储存方案}

观察 Euler 方程的守恒形式及 MacCormack 格式可知，时间推进的核心变量天然是
$\bm{Q}$，因此本程序以第一种方案为基础。同时，在内存允许的情况下，为原始量
$\bm{W}$ 也分配长期储存空间，但规定
\begin{equation}
\bm{Q}\text{ 为主状态},
\qquad
\bm{W}\text{ 为由 }\bm{Q}\text{ 同步得到的子状态}.
\end{equation}

$\bm{W}$ 不使用独立的数值格式推进，而是在 $\bm{Q}$ 更新后重新反解。一个完整时间步为
\begin{enumerate}[leftmargin=2em]
    \item 当前主状态 $\bm{Q}^{n}$ 已知；
    \item 由 $\bm{Q}^{n}$ 反解 $\bm{W}^{n}$；
    \item 用 $\bm{W}^{n}$ 计算声速、CFL、人工粘性和通量；
    \item 执行 MacCormack 预测步，得到 $\overline{\bm{Q}}$；
    \item 由 $\overline{\bm{Q}}$ 反解预测状态对应的原始量；
    \item 计算预测状态通量 $\overline{\bm{F}}$；
    \item 执行 MacCormack 校正步，得到 $\bm{Q}^{n+1}$；
    \item 由 $\bm{Q}^{n+1}$ 同步得到 $\bm{W}^{n+1}$；
    \item 输出当前状态并进入下一时间步。
\end{enumerate}

这种实现与仅临时计算 $\bm{W}$ 的第一种方案相比，主要区别是
$\bm{W}$ 拥有实际的数组储存空间，因此可直接查询、输出和调试。
代价是需要额外内存，并且必须保证 $\bm{W}$ 始终由最新的 $\bm{Q}$ 同步得到。

\subsection{时间变量储存策略}

\subsubsection{方案 A：全历史存储}

全历史存储为每一个时间层分配空间，保存
\begin{equation}
\bm{Q}^{0},\bm{Q}^{1},\ldots,\bm{Q}^{N_t}.
\end{equation}

该方案可以回看任意时刻、生成动画、进行完整时空分析、检查波传播过程，
并研究误差随时间的演化。但其内存和文件规模随网格数与时间步数的乘积增长，
计算时也会产生较大的缓存压力，因此不适合大型 CFD 计算。

例如，当
\begin{equation}
N_x=100000,\qquad
N_t=100000,\qquad
N_{\mathrm{var}}=3
\end{equation}
且每个 \texttt{double} 占用 $8$ 字节时，仅储存三项守恒量就需要
\begin{equation}
100000\times100000\times3\times8
=240\times10^9\ \text{bytes}
\approx240\ \text{GB}.
\end{equation}

该方案主要适用于小型算例、快速验证、后处理动画、程序调试和保存全部瞬态。

\subsubsection{方案 B：滚动时间层存储}

滚动存储只保存当前计算需要的少数时间层。对于本程序的 MacCormack 格式，
内存中主要保留
\begin{equation}
\bm{Q}^{n},
\qquad
\overline{\bm{Q}},
\qquad
\bm{Q}^{n+1}.
\end{equation}

完成一个时间步后，用 $\bm{Q}^{n+1}$ 覆盖旧的 $\bm{Q}^{n}$，再继续计算。
这种方式内存需求基本不随总时间步数增加，适合实际 CFD 时间推进。

\subsubsection{滚动存储与快照输出}

方案 B 仍可获得方案 A 的中间过程分析能力。本程序在内存中使用滚动存储，
同时每隔 $N$ 个时间步将当前流场写入磁盘快照文件：
\begin{equation}
\text{内存：滚动时间层}
\qquad
\text{磁盘：离散快照}.
\end{equation}

若未来需要研究某个中间状态，可以
\begin{enumerate}[leftmargin=2em]
    \item 找到对应时刻的 snapshot 文件；
    \item 读入其中保存的 $\bm{Q}$ 或 $\bm{W}$；
    \item 将该状态作为新的初始条件；
    \item 从该时间点继续计算。
\end{enumerate}

因此，本程序采用“内存滚动推进、磁盘按需保存快照”的策略，在控制内存消耗的同时，
保留对中间过程进行绘图、分析和重新计算的能力。


\subsection{空间与时间离散}

计算域离散为 $N_x$ 个节点：

\begin{equation}
x_i=x_{\min}+i\Delta x,
\qquad
\Delta x=\frac{x_{\max}-x_{\min}}{N_x-1}.
\end{equation}

当前程序并未采用自适应求解域（不是指自适应加密网格，而是指根据dx自适应数值计算域的空间。
主要原因是：固定计算域和均匀网格便于控制变量，使不同 $\Delta x$、CFL 数和人工粘性参数下的结果能够直接比较；

\subsection{CFL 时间步}

\begin{equation}
\Delta t
=
\mathrm{CFL}
\frac{\Delta x}
{\max_i(|u_i|+a_i)},
\qquad
a_i=\sqrt{\frac{\gamma p_i}{\rho_i}}.
\end{equation}

结合 Homework 3 中采用的冯诺依曼稳定性分析，可将离散格式对某一波数扰动的作用
写成放大因子 $G(k)$。在复平面上，稳定性要求任意波数 $k$ 对应的放大因子均满足
\begin{equation}
\left|G(k)\right|\leq 1,
\qquad \forall\,k,
\end{equation}
即所有波数扰动的放大因子轨迹必须位于单位圆内部或单位圆上。CFL 条件正是对
$\Delta t/\Delta x$ 的限制，使离散格式的放大因子不越出复平面上的单位圆，
从而保证任意傅里叶扰动不会随时间步进而无限放大。

\subsection{MacCormack 预测--校正格式}

前向通量差分预测、后向通量差分校正：

\begin{align}
\widetilde{\bm{Q}}_i^{\,n+1}
&=
\bm{Q}_i^n
-
\frac{\Delta t}{\Delta x}
\left(
\bm{F}_{i+1}^n-\bm{F}_i^n
\right),\\
\bm{Q}_i^{n+1}
&=
\frac{1}{2}
\left[
\bm{Q}_i^n+\widetilde{\bm{Q}}_i^{\,n+1}
-
\frac{\Delta t}{\Delta x}
\left(
\widetilde{\bm{F}}_i^{\,n+1}
-
\widetilde{\bm{F}}_{i-1}^{\,n+1}
\right)
\right].
\end{align}

关于格式阶数、差分方向配对方式以及在间断附近出现非物理振荡的原因，此处不赘述，在 文档2-结果与分析 中再作具体讨论

\subsection{人工粘性}

为比较不同物理量对间断位置的识别能力，程序允许在密度 $\rho$、速度 $u$
和压力 $p$ 之间自由选择人工粘性传感变量。令
\begin{equation}
\phi\in\{\rho,u,p\},
\end{equation}
则统一采用归一化二阶差分构造传感器：
\begin{equation}
\theta_i(\phi)=
\frac{
\left|\phi_{i+1}-2\phi_i+\phi_{i-1}\right|
}{
\left|\phi_{i+1}\right|
+2\left|\phi_i\right|
+\left|\phi_{i-1}\right|
+\varepsilon
},
\end{equation}
其中 $\varepsilon$ 是防止分母为零的小量。人工粘性强度为
\begin{equation}
\mu_i=\beta\,\theta_i(\phi),
\end{equation}
其中 $\beta$ 为人工粘性经验系数。

代码中三种选择共用同一套人工粘性算法，仅切换传入传感器的
$\phi$ 数组，因此可以在其他参数完全相同的条件下，直接比较
$\rho$、$u$ 和 $p$ 作为开关变量时的数值结果。

\subsection{边界条件}

当前程序使用零梯度外推：

\begin{equation}
\bm{Q}_0=\bm{Q}_1,
\qquad
\bm{Q}_{N_x-1}=\bm{Q}_{N_x-2}.
\end{equation}

说明这种边界条件适用的前提，以及计算域应足够大以避免波在终止时间前到达边界。

\section{程序结构与代码逻辑（伪代码形式呈现）}

\subsection{求解器结构体}

程序使用结构体统一保存：

\begin{itemize}[leftmargin=2em]
    \item 网格、时间和气体参数；
    \item 左右初始状态；
    \item 当前、预测和下一时间层守恒量；
    \item 当前与预测状态的通量；
    \item 原始量、声速及人工粘性传感器；
    \item 内存分配、初始化、推进、输出和释放等方法函数指针。
\end{itemize}

\begin{lstlisting}[caption={求解器结构体示意}]
typedef struct Riemann_1D_MacC_solver solver;

struct Riemann_1D_MacC_solver {
    int nx;
    double dx, dt, t, t_max, cfl, gamma;

    double *q1, *q2, *q3;
    double *q1_bar, *q2_bar, *q3_bar;
    double *q1_next, *q2_next, *q3_next;

    double *rho, *u, *p, *E, *a;

    int  (*allocate)(solver *self);
    void (*step_maccormack)(solver *self);
    void (*destroy)(solver *self);
};
\end{lstlisting}

\subsection{内存组织}

说明方案 B2.0：内存中只保留当前状态、预测状态和下一状态，计算完成后
用下一状态覆盖当前状态；中间过程通过定期写出快照文件保存。

\subsection{方法绑定与伪 OOP}

\begin{lstlisting}[caption={方法函数绑定示意}]
void solver_bind_methods(solver *self) {
    self->allocate = solver_allocate;
    self->step_maccormack = solver_step_maccormack;
    self->destroy = solver_destroy;
}
\end{lstlisting}

解释函数指针、\texttt{solver *self} 和
\texttt{self->member} 的含义。

\subsection{主程序流程}

\begin{enumerate}[leftmargin=2em]
    \item 创建默认求解器并绑定方法；
    \item 读取算例、网格和数值参数；
    \item 为所有数组分配堆内存；
    \item 初始化 Riemann 左右状态；
    \item 依据 CFL 计算时间步；
    \item 执行 MacCormack 预测与校正；
    \item 更新守恒量和原始量；
    \item 按间隔输出快照，最终输出结果；
    \item 释放动态内存。
\end{enumerate}

\section{关键实现细节}

\subsection{守恒量到原始量的转换}

\begin{align}
\rho &= q_1,\\
u &= \frac{q_2}{q_1},\\
p &= (\gamma-1)
\left(
q_3-\frac{q_2^2}{2q_1}
\right).
\end{align}

说明为何 $q_1>0$ 且 $p>0$ 是后续计算声速的必要条件。

\subsection{数值异常检测}

建议在最终程序中记录并检查：

\begin{itemize}[leftmargin=2em]
    \item $\rho\le 0$ 或 $p\le 0$；
    \item \texttt{NaN} 或 \texttt{Inf}；
    \item CFL 超出稳定范围；
    \item 时间步过小或达到最大迭代次数；
    \item 波到达计算域边界。
\end{itemize}

\subsection{快照与重启}

快照文件同时保存原始量和守恒量：

\begin{equation}
(x,\rho,u,p,E,q_1,q_2,q_3).
\end{equation}

说明快照用于绘图、动画和中间状态分析。若要实现严格 restart，
还应保存当前时间、步数、网格参数和气体参数，并增加读取函数。

\subsection{程序当前限制}

在此如实记录最终提交版本仍存在的限制，例如：

\begin{itemize}[leftmargin=2em]
    \item 人工粘性公式的经验性；
    \item MacCormack 对强间断的振荡；
    \item 边界条件较简单；
    \item 尚未实现高分辨率限制器；
    \item 尚未实现读取快照继续计算。
\end{itemize}

\section{本部分小结}

总结本报告介绍的数学模型、数值格式和代码架构。此处不展开数值结果，
结果、图像、误差与参数影响统一放入第二份报告。

\section*{参考资料}

\begin{enumerate}[leftmargin=2em]
    \item 填写课程教材或讲义。
    \item 填写 MacCormack 格式相关参考资料。
    \item 填写一维 Riemann 问题相关参考资料。
\end{enumerate}


\end{document}
